% Each \section is a new appendix (A, B, ...).
% Figures and tables are numbered within their appendix (A1, A2, ..., B1, ...),
% so put them in the appendix section they belong to.
% Put a \clearpage before each \section and at the end of the file,
% so each appendix's figures are flushed before the next appendix starts
% (otherwise they can float into the next appendix).
% \appendix and \noappendix are added by the build, don't add them here.
%
% Despite what the Copernicus template says, we don't use \appendixfigures or \appendixtables:
% they break the numbering.
% The template only asks for them if all figures and tables are put after the reference list
% (its "Option 2"). We put figures and tables in the sections they belong to (its "Option 1"),
% for which \appendix already numbers them per appendix (A1, A2, ..., B1, ...).
% \appendixfigures also steps the section counter, on top of the step from \section,
% so using it as well (tested 2026-10-07, one after each appendix \section)
% numbered the figures in Appendix B as C1, C2, ...

\clearpage
\section{Detailed methods figures}
\label{app:methods-figures}

\begin{figure}[hbt!]  % get latex to keep this figure after the section header, or at the bottom of this page or at the top of the next
\includegraphics[width=12cm]{<n2o-methods-appendix-figure>}
\caption{
    Working behind Figure \ref{fig:methods-n2o}.
    a), b) Demonstration of interpolation in the case where we have the most/least available values from the observation networks
    (but still pass our minimum data for interpolation requirement).
    c) Variance explained by each latitudinal gradient empirical orthogonal function (EOF).
    d) Latitudinal gradient principal components (PCs) derived from the observation network values.
}
\label{fig:methods-appendix-n2o}
\end{figure}

\begin{figure}[tp]
\includegraphics[width=12cm]{<co2-methods-appendix-figure>}
\caption{
    Working behind Figure \ref{fig:methods-co2}.
    a), b) Demonstration of interpolation in the case where we have the most/least available values from the observation networks
    (but still pass our minimum data for interpolation requirement).
    c) Variance explained by each seasonality change empirical orthogonal function (EOF).
    d) Seasonality change principal component (PC) derived from the observation network values.
    e) Regression between the seasonality change PC and the temperature - CO_2 concentration composite.
    f) Variance explained by each latitudinal gradient EOF.
    g) Latitudinal gradient PCs derived from the observation network values.
    h) Regression between the first latitudinal gradient PC and CO_2 emissions of geological origin.
}
\label{fig:methods-appendix-co2}
\end{figure}

\begin{figure}[tp]
\includegraphics[width=12cm]{<ch4-methods-appendix-figure>}
\caption{
    Working behind Figure \ref{fig:methods-ch4}.
    a), b) Demonstration of interpolation in the case where we have the most/least available values from the observation networks
    (but still pass our minimum data for interpolation requirement).
    c) Variance explained by each latitudinal gradient empirical orthogonal function (EOF).
    d) Latitudinal gradient principal components (PCs) derived from the observation network values.
    e) Regression between the first latitudinal gradient PC and CH_4 emissions of geological origin.
}
\label{fig:methods-appendix-ch4}
\end{figure}

% ZNTODO: check panel e) against the method
% (regression against total SSP2-4.5 emissions from RCMIP, see TODOs.md section 2)
\begin{figure}[tp]
\includegraphics[width=12cm]{<cfc12-methods-appendix-figure>}
\caption{
    Working behind Figure \ref{fig:methods-cfc12}.
    a), b) Demonstration of interpolation in the case where we have the most/least available values from the observation networks
    (but still pass our minimum data for interpolation requirement).
    c) Variance explained by each latitudinal gradient empirical orthogonal function (EOF).
    d) Latitudinal gradient principal components (PCs) derived from the observation network values.
    e) Regression between the first latitudinal gradient PC and CFC12 emissions.
}
\label{fig:methods-appendix-cfc12}
\end{figure}

\begin{figure}[tp]
\includegraphics[width=12cm]{<sf6-methods-figure>}
\caption{
    Production of the SF_6 datasets (for full details, see Section \ref{ssec:methods-cfc12-like}).
    a) Raw observation network data.
    b) Number of observations that feed into each of our 15\textdegree~latitudinal bins in each month.
    c) Location of each sampling site in the observation networks.
    d) Global-, annual-mean derived from the observation network values.
    e) Seasonality derived from the observation network values.
    f) Latitudinal gradient empirical orthogonal functions (EOFs) derived from the observation network values.
    g) Global-, annual-mean values for the entire output dataset timespan alongside the input data sources used.
    h) Latitudinal gradient principal components (PCs) for the entire output dataset timespan.
}
\label{fig:methods-sf6}
\end{figure}

\begin{figure}[tp]
\includegraphics[width=12cm]{<sf6-methods-appendix-figure>}
\caption{
    Working behind Figure \ref{fig:methods-sf6}.
    a), b) Demonstration of interpolation in the case where we have the most/least available values from the observation networks
    (but still pass our minimum data for interpolation requirement).
    c) Variance explained by each latitudinal gradient empirical orthogonal function (EOF).
    d) Latitudinal gradient principal components (PCs) derived from the observation network values.
    e) Regression between the first latitudinal gradient PC and SF_6 emissions.
}
\label{fig:methods-appendix-sf6}
\end{figure}

\clearpage
\section{Detailed results figures}
\label{app:results-figures}

% ZNTODO: the results figures for the other gases not shown in the main text (see TODOs.md, section 2)

\begin{figure}[hbt!]  % get latex to keep this figure after the section header, or at the bottom of this page or at the top of the next
\includegraphics[width=12cm]{<sf6-results-figure>}
\caption{
    SF_6 forcing dataset and comparison with other datasets.
    a) Surface plot of the last ten years of the native resolution dataset.
    b), c) Hemispheric- and global- monthly-mean values from 2000 to 2020 and 2020 onwards respectively.
    d), e), f) Hemispheric- and global-, annual-mean values from year 1 to 1750, 1750 to 2000 and 2000 onwards respectively.
    Note that each panel has its own vertical scale so the features of interest can be seen
    (but this also makes the differences in the pre-industrial period appear relatively bigger).
    In panels b) to f), the observation network values used as input data (Figure \ref{fig:methods-sf6}a)
    are shown for context, and coloured by latitude.
    g) The difference between the output CMIP7 dataset's global-, annual-mean and the CMIP6 global-, annual-mean.
    The right-hand axis shows the approximate difference in radiative forcing (for details, see Section \ref{ssec:results-radiative-forcing-comparisons-note}).
}
\label{fig:results-sf6}
\end{figure}

\begin{figure}[tp]
\includegraphics[width=12cm]{<cfc11eq-results-figure>}
\caption{
    CFC-11-eq forcing dataset.
    a) Surface plot of the last ten years of the native resolution dataset.
    b), c) Hemispheric- and global- monthly-mean values from 2000 to 2020 and 2020 onwards respectively.
    d), e), f) Hemispheric- and global-, annual-mean values from year 1 to 1750, 1750 to 2000 and 2000 onwards respectively.
    Note that each panel has its own vertical scale so the features of interest can be seen.
    IGCC has no CFC-11-eq, so there is no IGCC comparison.
    g) The difference between the output CMIP7 dataset's global-, annual-mean and the CMIP6 global-, annual-mean.
    The right-hand axis shows the approximate difference in radiative forcing (for details, see Section \ref{ssec:results-radiative-forcing-comparisons-note}).
}
\label{fig:results-cfc11eq}
\end{figure}

\clearpage
